\begin{lemma}
\label{editorial-source-lemma-integral-local}
Let $R \to S$ be a ring map.
Let $f_1, \ldots, f_n \in R$ generate the unit ideal.
\begin{enumerate}
\item If each $R_{f_i} \to S_{f_i}$ is integral, so is $R \to S$.
\item If each $R_{f_i} \to S_{f_i}$ is finite, so is $R \to S$.
\end{enumerate}
\end{lemma}

\begin{proof}
For (1), fix the original $s\in S$ and keep
$I=\{P\in R[x]:P^\varphi(s)=0\}$, the kernel of evaluation.
Let $J\subset R$ consist of zero and the leading coefficients
of the nonzero polynomials in $I$. This is an ideal. To add two
nonzero leading coefficients, multiply the corresponding polynomials
by powers of $x$ to make their degrees equal. If the sum of the
leading coefficients is nonzero, it is the leading coefficient of
the sum polynomial; if it is zero, it belongs to $J$ by definition.
For multiplication by $r\in R$, a nonzero product of leading
coefficients remains the leading coefficient of $rP$; a zero
product again belongs to $J$. Thus all ideal axioms hold.

For each original $f_i$, choose the local monic equation
$$
s^{d_i}+\sum_{j=0}^{d_i-1}
\varphi_{f_i}(a_{ij}/f_i^{c_{ij}})s^j=0
\quad\text{in }S_{f_i},
$$
where $d_i\geq1$, $a_{ij}\in R$ and $c_{ij}\geq0$.
Choose $N_i\geq c_{ij}$ for all $j$. The element
$\varphi(f_i^{N_i})s^{d_i}+
\sum_j\varphi(a_{ij}f_i^{N_i-c_{ij}})s^j$ becomes zero in
$S_{f_i}$, so some $M_i\geq0$ makes it zero upon multiplication
by $\varphi(f_i^{M_i})$. Consequently the original polynomial
$$
f_i^{N_i+M_i}x^{d_i}
+\sum_{j=0}^{d_i-1}a_{ij}f_i^{N_i+M_i-c_{ij}}x^j
$$
belongs to $I$. If $f_i^{N_i+M_i}\ne0$ it is its leading
coefficient; if it is zero it belongs to $J$ by definition.
In either case $f_i^{N_i+M_i}\in J$. A proper $J$ would lie
in a maximal ideal containing every $f_i$ (or would already
contain $1$ when an exponent is zero), contrary to the given
unit-ideal hypothesis. Therefore $1\in J$, so $I$ contains a
monic polynomial and $s$ is integral over $R$.

The maximal-ideal criterion above gives the same conclusion:
every maximal ideal omits some $f_i$ and the corresponding local
map is a further localization of $R_{f_i}\to S_{f_i}$.

For (2), choose finite $R_{f_i}$-module generators of $S_{f_i}$.
Write each as $s_{ij}/f_i^{a_{ij}}$, keeping its numerator
$s_{ij}\in S$ and exponent $a_{ij}\geq0$.
Let $N\subset S$ be the $R$-submodule generated by all these
finitely many numerators. Since $f_i$ is a unit after localization,
$N_{f_i}=S_{f_i}$. Thus the quotient $Q=S/N$ satisfies $Q_{f_i}=0$
for each $i$. For $q\in Q$ choose $b_i\geq1$ with $f_i^{b_i}q=0$.
The ideal generated by the $f_i^{b_i}$ is the unit ideal: if it were
proper, a maximal ideal containing it would contain all $f_i$,
contrary to the original unit-ideal assumption. Choose $c_i\in R$
with $\sum_i c_if_i^{b_i}=1$. Then
$q=\sum_i c_if_i^{b_i}q=0$. Hence $Q=0$ and the original numerators
$s_{ij}$ generate $S$ over $R$. The empty-cover case forces $R=0$
and hence $S=0$, where the assertion also holds.
\end{proof}